\section{An alternating two-parameter digamma identity}
\label{alt:sec:master}

The centered moments in the preceding sections use even polynomial weights.
An alternating outer sign has a complementary effect: odd centered weights
lead to odd-weight alternating double zeta values, which admit a classical
parity reduction. Summing these reductions over both even tail orders gives
a compact identity in digamma and trigonometric functions. The method of
turning products of tails into ordered sums is classical; see Hoffman
\cite{Hoffman}. The parity formula is due to the classical Euler-sum
theory and is proved in a generating-function form by Teo
\cite{Teo}; see also the descending-index formulation in
\cite[Eq.~(3.7)]{XuWang}. The contribution here is the resulting
two-parameter identity, its explicitly normalized odd-weight family, and
the rational-parameter square evaluations. No literature-priority claim
is made for individual special cases.

For $\sigma\in\{1,-1\}$ set
\begin{equation}\label{alt:eq:resolvent-tail}
 E_n^\sigma(x)=\sum_{k>n}\frac{\sigma^k}{k^2-x^2},
 \qquad x\notin\mathbb Z\setminus\{0\},\quad n\ge0.
\end{equation}
Each $E_n^\sigma$ is even in $x$. Its poles at integers of modulus at
most $n$ are absent. These functions are ordinary polygamma expressions:
\begin{align}
 E_n^+(x)&=\frac{\psi(n+1+x)-\psi(n+1-x)}{2x},
 \label{alt:eq:digamma-tail}\\
 E_n^-(x)&=\frac{(-1)^{n+1}}{2x}
       \{\mathcal B(n+1-x)-\mathcal B(n+1+x)\},
 \qquad
 \mathcal B(v)=\frac{\psi((v+1)/2)-\psi(v/2)}2.
 \notag
\end{align}
The apparent singularities at $x=0$ are removable. Formula
\eqref{alt:eq:digamma-tail} follows by partial fractions and the
convergent digamma difference series; the second formula follows from
$\mathcal B(v)=\sum_{j\ge0}(-1)^j/(j+v)$ for $\Re v>0$,
and then by continuation.

Define the even meromorphic functions
\begin{align}
 O_+(t)&=-\frac{\psi(1+t)+\psi(1-t)+2\gamma}{2},
 \label{alt:eq:O}\\
 O_-(t)&=-\log2+O_+(t/2)-O_+(t),\notag
\end{align}
and put $R_+(t)=\cot(\pi t)$ and $R_-(t)=\csc(\pi t)$.
For $xy\ne0$, abbreviate
\begin{equation}\label{alt:eq:UV}
 U_\sigma(x,y)=\frac{O_\sigma(x+y)-O_\sigma(x-y)}{2xy},
 \qquad
 V_\sigma(x,y)=\frac{O_\sigma(x+y)+O_\sigma(x-y)}2.
\end{equation}

\begin{theorem}[The two-parameter alternating identity]
\label{alt:thm:master}
Let $\alpha,\beta\in\{1,-1\}$ and $c=-\alpha\beta$. Initially assume
$|x|+|y|<1$ and $xy\ne0$. Then
\begin{align}
 &\mathcal K_{\alpha,\beta}(x,y)
 :=\sum_{n=0}^\infty(-1)^n(2n+1)E_n^\alpha(x)E_n^\beta(y)
 \notag\\
 &=\frac\pi2\bigg\{
 \frac{R_\alpha(x)}xO_{-\beta}(y)
 +\frac{R_\beta(y)}yO_{-\alpha}(x)
 -yR_c(y)U_\alpha(x,y)-xR_c(x)U_\beta(x,y)
 \notag\\[-2pt]
 &\hspace{37mm}
 -\frac{R_c(x)}xV_{-\beta}(x,y)
 -\frac{R_c(y)}yV_{-\alpha}(x,y)\bigg\}.
 \label{alt:eq:master}
\end{align}
The series converges absolutely and locally normally if at least one of
$\alpha,\beta$ is $-1$. For $\alpha=\beta=1$ it converges ordinarily
and locally uniformly; after subtracting each summand's value at
$(0,0)$ the remaining series is absolutely locally normally
convergent. Both sides extend meromorphically to $\mathbb C^2$.
Possible genuine poles occur only when $x$ or $y$ is a nonzero integer.
In particular the apparent singularities at $xy=0$, and those on
$x\pm y\in\mathbb Z$ away from the coordinate poles, are removable.
The identity is even separately in $x$ and $y$ and is unchanged by
simultaneously exchanging $(\alpha,x)$ with $(\beta,y)$.
\end{theorem}

\subsection{An ordinary finite-cutoff identity}

For $s\ge2$, write
\[
 T_s^\sigma(n)=\sum_{k>n}\frac{\sigma^k}{k^s},\qquad
 Z(r,s;\sigma,\tau)
 =\sum_{k>m\ge1}\frac{\sigma^k\tau^m}{k^r m^s}.
\]
The convention for $Z$ is the decreasing-index convention. Put
$\mathsf L_s(\sigma)=\operatorname{Li}_s(\sigma)$ for $s\ge2$ and
$\mathsf L_1(-1)=-\log2$.

\begin{lemma}[Odd alternating weights]
\label{alt:lem:finitecutoff}
Let $p,q\ge2$ be integers and let $d$ be odd with
$1\le d<\min(p,q)$. Then
\begin{align}
 &\sum_{n\ge0}(-1)^n\{n^d+(n+1)^d\}
                    T_p^\alpha(n)T_q^\beta(n)\notag\\
 &\quad=-Z(p,q-d;\alpha,-\beta)
        -Z(q,p-d;\beta,-\alpha)
        -\mathsf L_{p+q-d}(-\alpha\beta).
 \label{alt:eq:weighted-tail}
\end{align}
The left side is absolutely convergent except possibly when
$(p,q,d,\alpha,\beta)=(2,2,1,1,1)$; in that case it converges as
an ordinary alternating series. Every double series on the right is
absolutely convergent.
\end{lemma}

\begin{proof}
For $m\ge1$, cancellation of consecutive powers gives
\[
 \sum_{n=0}^{m-1}(-1)^n\{n^d+(n+1)^d\}
       =(-1)^{m-1}m^d.
\]
First sum over the finite range $0\le n\le N$. For each fixed $n$
the two tail series are absolutely convergent, so their product may
be expanded and the finite $n$-sum interchanged. The inner weight
becomes $(-1)^{\min(k,m,N+1)-1}\min(k,m,N+1)^d$.
The limiting double series is dominated by
$\min(k,m)^d/(k^p m^q)$, whose sum is finite because $q-d,p-d\ge1$
and $p,q\ge2$. For example the region $k>m$ is
\[
 \sum_{m\ge1}m^{d-q}\sum_{k>m}k^{-p}
       =O\left(\sum_{m\ge1}m^{d+1-p-q}\right)<\infty.
\]
Dominated convergence is therefore applicable to the double series.
Separating $k>m$, $m>k$, and $k=m$ proves
\eqref{alt:eq:weighted-tail}. Equivalently, the additional term in the
finite-cutoff formula when both indices exceed $N+1$ has modulus at
most $(N+1)^d T_p^+(N)T_q^+(N)=O(N^{d+2-p-q})$, which tends to zero.

For absolute convergence of the original series use
$T_s^+(n)=O(n^{1-s})$ and $T_s^-(n)=O(n^{-s})$.
Under the stated range the only unresolved exponent is $-1$ in the
listed case. There
$(2n+1)T_2^+(n)^2=2/n+O(n^{-2})$ for $n\ge1$, so an ordinary
alternating harmonic series plus an absolutely convergent series
proves convergence. This also fixes the boundary value without
assigning a value to a divergent series.
\end{proof}

\subsection{The classical parity input, with its endpoint convention}

The following formula is the alternating double-Euler parity theorem
\cite{Teo,XuWang}. We record it in the convention used here.
For $r\ge2$, $s\ge1$, and $w=r+s$ odd, set
\[
 \mu_j=(-1)^r\left\{
 \binom{j-1}{r-1}\mathsf L_j(\sigma)
 +\binom{j-1}{s-1}\mathsf L_j(\tau)\right\}.
\]
Binomial coefficients outside $0\le k\le n$ are zero. Then
\begin{align}
 Z(r,s;\sigma,\tau)
 ={}&\frac12\left\{-\mathsf L_w(\sigma\tau)
 +(1+(-1)^r)\mathsf L_r(\sigma)\mathsf L_s(\tau)+\mu_w\right\}
 \notag\\
 &-\sum_{k=1}^{(w-1)/2}\mathsf L_{2k}(\sigma\tau)\mu_{w-2k}.
 \label{alt:eq:parity}
\end{align}
The exceptional notation $\mathsf L_1(1)$ is interpreted as one formal
symbol $T$ while this finite expression is assembled. It cancels before
the expression is evaluated. Indeed it can occur only when $s=1$ and
$\tau=1$; then $r$ is even, and the first product contributes
$\mathsf L_r(\sigma)T$, while the last summand contributes its negative.
There is no remaining regularization parameter. Thus one may set
$\mathsf L_1(1)=0$ in a coefficient implementation after this cancellation
has been established. The value zero is a bookkeeping convention, not
an asserted value of the divergent harmonic series.

Combining \eqref{alt:eq:weighted-tail} and \eqref{alt:eq:parity}
already proves an explicit all-order family whenever $p+q$ is even.
Only single zeta values and $\log2$ remain, since
$\mathsf L_s(-1)=(2^{1-s}-1)\zeta(s)$ for $s\ge2$.
No period-independence assumption is involved.
The canonical weights exhaust the corresponding odd polynomial space:
for odd $d$,
$P_d(X)=(X-1/2)^d+(X+1/2)^d$ is odd with leading term $2X^d$.
Consequently the $P_d$ with $d<\min(p,q)$ form a triangular basis
for all odd centered polynomial weights in that degree range.
For $p+q$ even, every such weighted moment therefore has the same
finite reduction to single zeta values and $\log2$.

\subsection{Proof of the two-parameter identity}

Put
\[
 A_\sigma(t)=\sum_{a\ge1}\mathsf L_{2a}(\sigma)t^{2a-2}
       =\frac1{2t^2}-\frac{\pi R_\sigma(t)}{2t}.
\]
The equality follows from the cotangent partial-fraction expansion and
its alternating counterpart. Moreover, for $|t|<1$,
\[
 O_\sigma(t)=\sum_{j\ge0}\mathsf L_{2j+1}(\sigma)t^{2j},
\]
with the already justified convention $\mathsf L_1(1)=0$.
The familiar digamma Taylor expansion gives \eqref{alt:eq:O}.

We first verify the analytic legitimacy of the coefficient calculation.
Uniformly for $x$ in a compact subset avoiding the coordinate poles,
\[
 E_n^+(x)=O(n^{-1}),\quad E_n^-(x)=O(n^{-2}),\quad
 E_n^\sigma(x)-E_n^\sigma(0)=O(n^{-3}).
\]
The last estimate follows directly by subtracting the two rational
summands in \eqref{alt:eq:resolvent-tail}. Consequently, if one sign is
negative, the defining series for $\mathcal K$ converges absolutely
normally. In the positive-positive case its difference from the series
at $(0,0)$ is absolutely normally convergent, with summands $O(n^{-3})$.
The scalar series at $(0,0)$ converges by Lemma~\ref{alt:lem:finitecutoff}.
Cauchy's coefficient formula therefore gives
\begin{equation}\label{alt:eq:master-coeff}
 [x^{2a-2}y^{2b-2}]\mathcal K_{\alpha,\beta}
 =-Z(2a,2b-1;\alpha,-\beta)
  -Z(2b,2a-1;\beta,-\alpha)
  -\mathsf L_{2a+2b-1}(c).
\end{equation}

Insert \eqref{alt:eq:parity}. Its two terms
$-\mathsf L_{2a+2b-1}(c)/2$ cancel the diagonal term in
\eqref{alt:eq:master-coeff}. Summation of the remaining binomial terms
uses only the even and odd parts of the binomial theorem. For $N\ge1$,
\begin{align*}
 \sum_{a=1}^{N}\binom{2N}{2a-1}x^{2a-2}y^{2N-2a}
 &=\frac{(x+y)^{2N}-(x-y)^{2N}}{2xy},\\
 \sum_{a=0}^{N}\binom{2N}{2a}x^{2a}y^{2N-2a}
 &=\frac{(x+y)^{2N}+(x-y)^{2N}}2.
\end{align*}
For completeness, the complete expression before replacing $A_\sigma$
by trigonometric functions is
\begin{align}
 \mathcal K_{\alpha,\beta}
 ={}&\left(y^2A_c(y)-\tfrac12\right)U_\alpha
    +\left(x^2A_c(x)-\tfrac12\right)U_\beta\notag\\
 &+\left(A_c(x)-\tfrac1{2x^2}\right)V_{-\beta}
    +\left(A_c(y)-\tfrac1{2y^2}\right)V_{-\alpha}\notag\\
 &+\left(\tfrac1{2x^2}-A_\alpha(x)\right)O_{-\beta}(y)
    +\left(\tfrac1{2y^2}-A_\beta(y)\right)O_{-\alpha}(x).
 \label{alt:eq:master-before-trig}
\end{align}
All the power series used here converge absolutely when
$|x|+|y|<1$. Substitution of $A_\sigma$ proves
\eqref{alt:eq:master}. If the formal constant $T$ is retained in $O_+$,
it cancels directly in this last expression: when $-\beta=1$ one has
$c=\alpha$, and when $-\alpha=1$ one has $c=\beta$.

The left side is holomorphic on
$(\mathbb C\setminus(\mathbb Z\setminus\{0\}))^2$ by the same normal
convergence argument; near a coordinate pole only finitely many tail
terms are singular. The right side is meromorphic on $\mathbb C^2$.
The meromorphic identity theorem extends their equality from the
initial open set. This proves all the removability and continuation
claims in Theorem~\ref{alt:thm:master}. Evenness and exchange symmetry
follow directly from the left side.

There is also a direct check of the coordinate axes. Pair the first
and fifth terms in the braces of \eqref{alt:eq:master} as
\[
 \frac{R_\alpha(x)-R_c(x)}xO_{-\beta}(y)
 +\frac{R_c(x)}x\{O_{-\beta}(y)-V_{-\beta}(x,y)\}.
\]
All $R_\sigma$ have the same principal part $1/(\pi x)$, so the
first quotient is holomorphic at zero. The second difference is
divisible by $x^2$, so the second product is also holomorphic.
The symmetric pair handles $y=0$. Finally, the difference of any
even power evaluated at $x+y$ and $x-y$ is divisible by $xy$,
which proves that $U_\sigma$ is holomorphic on both axes.

\subsection{A diagonal square with elementary rational-parameter values}

\begin{corollary}[A closed alternating square]
\label{alt:cor:square}
For every $x\notin\mathbb Z\setminus\{0\}$,
\begin{equation}\label{alt:eq:diagonal-square}
 \boxed{\displaystyle
 \sum_{n\ge0}(-1)^n(2n+1)
       \left(\sum_{k>n}\frac{(-1)^k}{k^2-x^2}\right)^2
 =-\frac{\pi\{\psi(1+x)+\psi(1-x)+2\gamma\}}
             {4x\sin(\pi x)}.}
\end{equation}
The series is absolutely locally normally convergent. Its value at
$x=0$ is $\zeta(3)/2$.
\end{corollary}
\begin{proof}
Set $\alpha=\beta=-1$ and $y=x$ in
\eqref{alt:eq:master}. Here $c=-1$, $O_+(0)=0$, and
\[
 O_-(2x)-O_-(0)=O_+(x)-O_+(2x).
\]
The two terms at $2x$ cancel, leaving
$\pi O_+(x)/(2x\sin\pi x)$, which is
\eqref{alt:eq:diagonal-square}. At zero,
$O_+(x)=\zeta(3)x^2+O(x^4)$ and
$x\sin(\pi x)=\pi x^2+O(x^4)$. This proves the value and the
removable interpretation.
\end{proof}

The complete diagonal can be written as one symmetric matrix. With
$c_x=\cos(\pi x)$ and the signs ordered as $+,-$,
\begin{equation}\label{alt:eq:diagonal-matrix}
 \bigl(\mathcal K_{\alpha,\beta}(x,x)\bigr)_{\alpha,\beta\in\{+,-\}}
 =\frac{\pi}{2x\sin(\pi x)}
 \begin{pmatrix}
 2c_xO_-(x)-O_+(x)+2\log2 & O_-(x)+c_x\log2\\
 O_-(x)+c_x\log2 & O_+(x)
 \end{pmatrix}.
\end{equation}
This follows from the same substitution $y=x$ and
$O_+(2x)+O_-(2x)=O_+(x)-\log2$ in the other sign cases.
In particular,
\begin{equation}\label{alt:eq:half-matrix}
 \bigl(\mathcal K_{\alpha,\beta}(1/2,1/2)\bigr)_{\alpha,\beta}
 =\pi\begin{pmatrix}1&-1\\-1&2\log2-1\end{pmatrix}.
\end{equation}
Here $E_n^+(1/2)=1/(n+1/2)$, so the positive-positive entry also
reduces directly to the ordinary alternating arctangent series for
$\pi$. This example lies on the apparent oblique polar line $x+y=1$
of the unsimplified master formula, and makes its removability explicit.

For $x=r/N$, $1\le r<N$, Gauss's digamma formula
\cite[Eq.~(5.4.19)]{DLMFGamma} gives the explicit
finite expression
\begin{equation}\label{alt:eq:rational-O}
 O_+(r/N)=\log(2N)-\frac{N}{2r}
 -2\sum_{j=1}^{\lfloor(N-1)/2\rfloor}
       \cos\frac{2\pi rj}{N}\log\sin\frac{\pi j}{N}.
\end{equation}
Thus every such square is $\pi$ times an explicit linear combination,
with algebraic coefficients, of $1$ and logarithms of positive
algebraic numbers. For example,
\begin{align}
 \sum_{n\ge0}(-1)^n(2n+1)
  \left(\sum_{k>n}\frac{(-1)^k}{k^2-1/4}\right)^2
 &=\pi(2\log2-1),\label{alt:eq:half-example}\\
 \sum_{n\ge0}(-1)^n(2n+1)
  \left(\sum_{k>n}\frac{(-1)^k}{k^2-1/9}\right)^2
 &=\frac{3\sqrt3\pi}{2}(\log3-1),\label{alt:eq:third-example}\\
 \sum_{n\ge0}(-1)^n(2n+1)
  \left(\sum_{k>n}\frac{(-1)^k}{k^2-1/16}\right)^2
 &=2\sqrt2\pi(3\log2-2).\label{alt:eq:quarter-example}
\end{align}
The squares are algebraic squares throughout; no complex conjugation
is inserted when $x$ is complex.

Expanding \eqref{alt:eq:diagonal-square} produces a short all-order
convolution. With $\eta(s)=(1-2^{1-s})\zeta(s)$ for $s>1$, for every
$m\ge2$ one has
\begin{equation}\label{alt:eq:convolution}
 \sum_{a=1}^{m-1}\sum_{n\ge0}(-1)^n(2n+1)
        T_{2a}^-(n)T_{2m-2a}^-(n)
 =\frac12\zeta(2m-1)
   +\sum_{j=1}^{m-2}\eta(2j)\zeta(2m-2j-1).
\end{equation}
Indeed
$\pi x/\sin\pi x=1+2\sum_{j\ge1}\eta(2j)x^{2j}$;
coefficient extraction is legitimate by absolute normal convergence.

\begin{corollary}[An anchored Gamma primitive]
\label{alt:cor:gamma-primitive}
For $|u|<1$,
\begin{equation}\label{alt:eq:gamma-primitive}
 \int_0^u x\sin(\pi x)\mathcal K_{-,-}(x,x)\,dx
 =-\frac\pi4\left\{
 \log\Gamma(1+u)-\log\Gamma(1-u)+2\gamma u\right\}.
\end{equation}
The logarithms are their analytic branches on the unit disk,
normalized by $\log\Gamma(1)=0$. The path of integration may be any
path in that disk, and the defining series for $\mathcal K_{-,-}$
may be integrated term by term. In particular the value at $u=1/2$
is $\pi(\log2-\gamma)/4$.
\end{corollary}
\begin{proof}
Multiply \eqref{alt:eq:diagonal-square} by $x\sin(\pi x)$.
Its right side is the derivative of the expression in braces
in \eqref{alt:eq:gamma-primitive}, multiplied by $-\pi/4$.
The anchored logarithms exist because $\Gamma(1\pm u)$ are
holomorphic and nonzero on the simply connected unit disk.
Both sides vanish at zero. Absolute local normal convergence,
already proved, justifies exchanging the path integral and sum.
Finally $\Gamma(3/2)/\Gamma(1/2)=1/2$ gives the stated value.
\end{proof}

Some individual coefficients of the full two-variable theorem are
also useful. Let
$C_{p,q}^{\alpha,\beta}=\sum_{n\ge0}(-1)^n(2n+1)
T_p^\alpha(n)T_q^\beta(n)$. Then
\begin{align}
 C_{2,2}^{+,+}&=\frac{\pi^2\log2}{2}-\frac54\zeta(3),
 & C_{2,2}^{+,-}&=-\frac{\pi^2\log2}{4}-\frac38\zeta(3),
 \label{alt:eq:coeff22}\\
 C_{2,2}^{-,-}&=\frac12\zeta(3),
 &C_{2,4}^{+,+}&=\frac{\pi^4\log2}{48}
                   +\frac{\pi^2\zeta(3)}{12}-\frac{23}{32}\zeta(5),
 \label{alt:eq:coeff24}\\
 C_{4,4}^{+,+}&=-\frac{335}{64}\zeta(7)
                 +\frac{13\pi^2}{48}\zeta(5)
                 +\frac{\pi^4}{32}\zeta(3).
 \label{alt:eq:coeff44}
\end{align}
Only $C_{2,2}^{+,+}$ is conditionally convergent. As an example with a
higher odd centered weight, \eqref{alt:eq:weighted-tail} gives the
absolutely convergent identity
\begin{equation}\label{alt:eq:cubicweight}
 \sum_{n\ge0}(-1)^n\{n^3+(n+1)^3\}T_4^+(n)^2
 =\frac{\pi^4\log2}{24}+\frac{\pi^2\zeta(3)}8
       -\frac{49}{16}\zeta(5).
\end{equation}

\subsection{All repeated rational resolvents}

The same identity supplies explicit derivatives of every order. Define
\[
 E_{n,r}^{\sigma}(x)=\sum_{k>n}
       \frac{\sigma^k}{(k^2-x^2)^{r+1}},\qquad r\ge0,
\]
and let $\partial_{x^2}=(2x)^{-1}\partial_x$ on even analytic
functions, with its removable meaning at $x=0$.

\begin{corollary}[Repeated-resolvent products]\label{alt:cor:derivatives}
For integers $r,s\ge0$ and parameters outside the coordinate poles,
\begin{equation}\label{alt:eq:derivatives}
 \sum_{n\ge0}(-1)^n(2n+1)
        E_{n,r}^{\alpha}(x)E_{n,s}^{\beta}(y)
 =\frac{1}{r!s!}
      \partial_{x^2}^{\,r}\partial_{y^2}^{\,s}
                         \mathcal K_{\alpha,\beta}(x,y).
\end{equation}
The right side is an explicit finite expression obtained by
differentiating \eqref{alt:eq:master}. Whenever $r+s>0$, the series
is absolutely locally normally convergent for all sign pairs.
\end{corollary}
\begin{proof}
Termwise differentiation of the individual rational tails gives
$\partial_{x^2}^{\,r}E_n^\alpha(x)=r!E_{n,r}^\alpha(x)$.
For compact parameter sets avoiding poles, the positive-sign estimate
is $E_{n,r}^+(x)=O(n^{-2r-1})$; an alternating sign improves it to
$O(n^{-2r-2})$. Thus a derivative of positive total order makes the
outer series absolutely summable, even when both signs are positive.
The locally uniform identity in Theorem~\ref{alt:thm:master} and
Cauchy's differentiation formula justify the operation. Evenness
gives the stated removable interpretation at the axes.
\end{proof}

\subsection{A colored polylogarithmic function and its anchored primitive}

The same finite-difference mechanism has an interior-variable form that
also fixes an antiderivative. For $p,q\ge2$, $|a|,|b|\le1$, and
$|z|<1$, put
$T_p^a(n)=\operatorname{Li}_p(a)-\sum_{k=1}^n a^k/k^p$.
Use the colored convention
\[
 \operatorname{Li}_{s_1,\ldots,s_d}(z_1,\ldots,z_d)
 =\sum_{n_1>\cdots>n_d\ge1}
       \frac{z_1^{n_1}\cdots z_d^{n_d}}
            {n_1^{s_1}\cdots n_d^{s_d}}.
\]

\begin{proposition}[Colored generating function and primitive]
\label{alt:prop:GF}
Writing $F(z)=\sum_{n\ge0}T_p^a(n)T_q^b(n)z^n$, one has
\begin{align}
 (1-z)F(z)={}&\operatorname{Li}_p(a)\operatorname{Li}_q(b)
 -\operatorname{Li}_q(b)\operatorname{Li}_p(az)
 -\operatorname{Li}_p(a)\operatorname{Li}_q(bz)\notag\\
 &+\operatorname{Li}_{p,q}(az,b)
  +\operatorname{Li}_{q,p}(bz,a)
  +\operatorname{Li}_{p+q}(abz).
 \label{alt:eq:GF}
\end{align}
The primitive normalized to vanish at zero is
\begin{align}
 \int_0^zF(t)\,dt={}&
  \operatorname{Li}_p(a)\operatorname{Li}_q(b)\operatorname{Li}_1(z)
 -\operatorname{Li}_q(b)\operatorname{Li}_{1,p}(z,a)
 -\operatorname{Li}_p(a)\operatorname{Li}_{1,q}(z,b)\notag\\
 &+\operatorname{Li}_{1,p,q}(z,a,b)
  +\operatorname{Li}_{1,q,p}(z,b,a)
  +\operatorname{Li}_{1,p+q}(z,ab).
 \label{alt:eq:primitive}
\end{align}
The integral is along any path in the unit disk. Every function in the
formulas is defined by an absolutely locally normally convergent series.
\end{proposition}
\begin{proof}
The identity
\[
 T_p^a(n-1)T_q^b(n-1)-T_p^a(n)T_q^b(n)
 =\frac{a^n}{n^p}T_q^b(n)+\frac{b^n}{n^q}T_p^a(n)
   +\frac{(ab)^n}{n^{p+q}}
\]
proves \eqref{alt:eq:GF} after multiplication by $z^n$ and summation.
All rearrangements are absolute in the disk. The general derivative
\[
 \frac{d}{dz}\operatorname{Li}_{1,s_1,\ldots,s_d}
                (z,c_1,\ldots,c_d)
 =\frac{\operatorname{Li}_{s_1,\ldots,s_d}
                (zc_1,c_2,\ldots,c_d)}{1-z}
\]
follows by summing the largest index first. It proves the displayed
primitive term by term. Each term vanishes at $z=0$, which fixes its
additive constant and branch.
\end{proof}
The generating function and primitive are finite-difference consequences,
not assertions of general depth reduction. The additional reduction in
Theorem~\ref{alt:thm:master} uses the specific alternating weight and
odd-weight parity input. It does not imply the outstanding Gaussian
$S_6$ or revised $S_8$ identities.
